\section{Análisis de Resultados}\label{sec:analisis}

Los resultados obtenidos indican que el caudal de diseño de 2 260 CFM satisface el requerimiento de 12 renovaciones de aire por hora para el volumen de 320 m$^3$. Este valor se encuentra dentro de los rangos recomendados por ASHRAE 170 para espacios de laboratorio y por el WHO Laboratory Biosafety Manual para instalaciones BSL-2, con margen de seguridad respecto al mínimo de 6 ACH.

La potencia teórica del ventilador, calculada en 0.320 kW para un escenario de presión total de 165 Pa a las condiciones del sitio (Cajicá, 2 558 msnm, $\rho$ = 0.88 kg/m$^3$), es consistente con el orden de magnitud esperado para un ventilador axial de impulsión de caudal medio en aplicaciones de ventilación de laboratorio. La selección provisional de un motor de 0.75 HP TEFC anticorrosivo proporciona un factor de seguridad adecuado para compensar pérdidas adicionales por ensuciamiento de filtros, variaciones de eficiencia y condiciones de operación no ideales; el valor definitivo se confirmará con la curva de catálogo del axial seleccionado.

La velocidad de impulsión de 8 m/s en la boca del ventilador se ubica dentro del rango recomendado de 6 a 12 m/s para ventiladores centrífugos y axiales de impulsión directa. Esta velocidad permite un ingreso de aire ordenado sin generar ruido excesivo ni corrientes locales molestas en el interior del recinto.

La velocidad de descarga de 3.0 m/s en las rejillas de salida responde a un criterio de confort y de pérdida de presión baja en la descarga libre: una velocidad excesiva aumentaría el ruido y la velocidad del aire en las zonas cercanas a las rejillas. El valor seleccionado se encuentra dentro del rango aceptable de 2.5 a 4.0 m/s para aplicaciones de laboratorio y representa una pérdida de descarga de solo 11 Pa a la densidad del sitio.

La configuración de descarga adoptada, definida por BRINSA, dispone las tres rejillas apiladas verticalmente, una sobre otra, en la pared opuesta a la inyección. Esta disposición concentra las salidas en una sola columna vertical y fue la configuración evaluada en el modelo CFD; como se discute a continuación, la simulación confirma que el caudal se distribuye entre los tres niveles y que no se generan zonas de estancamiento severo.

El modelo CFD, resuelto con salidas tipo \texttt{pressure outlet} a \SI{0}{Pa} gauge (descarga libre a la atmósfera), confirma el comportamiento esperado del sistema de ventilación. Las Figuras \ref{fig:cfd_streamlines} y \ref{fig:cfd_streamlines_b} muestran que el aire inyectado describe un chorro dirigido hacia la pared de descarga, alineado a altura media con la rejilla central de la columna; al aproximarse a la pared opuesta, las líneas de corriente se abren verticalmente para alimentar las rejillas superior e inferior, sin retorno hacia la inyección, lo cual evidencia una distribución de flujo ordenada a través del recinto. La escala de velocidad alcanza aproximadamente \SI{8}{m/s} en la boca de impulsión y desciende progresivamente hacia las salidas.

Las Figuras \ref{fig:cfd_vectores_3d} y \ref{fig:cfd_vectores_3d_b} evidencian en el corte vertical la difusión del chorro y la formación de celdas de recirculación en los tercios superior e inferior del recinto. Aunque existe recirculación, los vectores mantienen dirección de salida en las tres rejillas apiladas, sin indicios de flujo entrante. Esta recirculación es aceptable en un sistema de ventilación por dilución siempre que no genere zonas de estancamiento prolongado; la propia recirculación es el mecanismo que alimenta las rejillas de los extremos de la columna, por lo que contribuye a distribuir el caudal de descarga entre los tres niveles.

Las Figuras \ref{fig:cfd_planta_horizontal} y \ref{fig:cfd_planta_horizontal_b}, correspondientes al corte horizontal a la altura de la rejilla central, confirman el barrido longitudinal entre la pared de inyección y la pared de descarga. Los vectores de mayor módulo se concentran en el eje del chorro, mientras que las zonas periféricas presentan dos vórtices de arrastre simétricos de baja velocidad. Este patrón es típico de una inyección localizada sin ductos de distribución; la magnitud de los vórtices no compromete el objetivo de renovación porque el caudal total (12 ACH) garantiza el reemplazo del aire.

Las Figuras \ref{fig:cfd_vectores_corte} y \ref{fig:cfd_vectores_corte_b} refuerzan las observaciones anteriores en el plano vertical longitudinal que atraviesa la columna de rejillas: el chorro de inyección mantiene su dirección hasta la proximidad de la pared de descarga, alimenta directamente la rejilla central y, mediante las celdas de recirculación superior e inferior, distribuye el resto del caudal hacia las rejillas de los extremos. Las tres rejillas descargan con velocidades faciales del orden de \SI{3}{m/s}, coherentes con el dimensionamiento, y no se observan corrientes de retorno significativas hacia la inyección ni fugas no modeladas. La distribución de velocidad es coherente con el balance de masas teórico y con el objetivo de ventilación por dilución del recinto.

Finalmente, el balance de masas entre la entrada y la salida del recinto se cierra con una discrepancia inferior al 0.1 \%, lo cual valida la coherencia dimensional y cinemática de los datos de contorno propuestos para el CFD.
